% ----------------------------------------------------------
% Apêndices
% ----------------------------------------------------------

% ---
% Inicia os Apêndices
% ---
\begin{apendicesenv}

% Imprime uma página indicando o início dos Apêndices
\partapendices

% ----------------------------------------------------------
\chapter{Vídeos demonstrando o rastreio pelo sistema simulado no software MuJoCo}\label{cap:apendice_umbra}
% ----------------------------------------------------------

Para exemplificar o funcionamento do protótipo, foi desenvolvido um modelo simulado no software MuJoCo (discutido no capítulo \ref{sec:MuJoCo}).
O primeiro vídeo pode ser encontrado na referência \cite{youtube_umbra} e na Figura \ref{fig:QR_Umbra} e contém uma demonstração do sistema rastreando o satélite Umbra-07.

\begin{figure}[htbp]
    \centering
    \includegraphics[width=0.5\textwidth]{figs/QRCODE_passagem_umbra.png}
    \caption{Código QR do vídeo do sistema rastreando o satélite Umbra-07}
    \legend{Fonte: Autor(2026)}
    \label{fig:QR_Umbra}
\end{figure}

Em 3:55, é possível observar o aumento da velocidade da passagem quando o satélite se aproxima de elevações mais altas.
O Umbra-07 atinge a sua maior elevação em 5:02, quando inicia a reduzir a elevação em relação à antena.
A passagem é finalizada em 10:57, no qual a elevação do satélite chega a zero e o sistema retorna à sua posição de origem.

\begin{figure}[htbp]
    \centering
    \includegraphics[width=0.5\textwidth]{figs/QRCODE_troca_satelites.png}
    \caption{Código QR do vídeo do sistema rastreando o satélite Umbra-07}
    \legend{Fonte: Autor(2026)}
    \label{fig:QR_Troca_Sats}
\end{figure}

O segundo vídeo é referente ao filtro lidando com troca de satélites alvo, presente na referência \cite{youtube_troca_satelite} e no Código QR da Figura \ref{fig:QR_Troca_Sats}.
O vídeo se inicia com o usuário selecionando o satélite COSMOS 2561 a ser rastreado utilizando o SatDump.
Em 0:06, o sistema é acionado e a antena é alinhada com o satélite, aproveitando a primera volta do sistema.
Em 0:34, o satélite KUPIER-00183 é selecionado e a antena inicia o novo rastreio.
O envio do rastreio é desabilitado em 1:20, e é possível perceber o sistema divergindo da posição real.

Em 1:51, o sistema para de rastrear o ONEWEB-0399 e começa a rastrear o satélite YAOGAN-30 05C, sendo possível verificar o sistema oscilando para se corrigir à nova órbita.
O sistema se estabiliza por volta de 2:12, conseguindo rastrear a nova trajetória.
Em 2:46, o algoritmo verifica que não está recebendo mais informações do satélite e retorna à sua posição inicial, mas logo inicia o rastreio do KSM1-C.
A troca do rastreio do KSM1-C para o 2024-205G em 3:05 provoca uma outra oscilação no sistema enquanto o filtro captura as informações da nova órbita.

A alteração do ONEWEB-0095 para o CE-SAT-IE em 3:40 demonstra que o filtro não sofre muitas oscilações quando as trajetórias são próximas.
Em 4:55, há uma demonstração do filtro quando este recebe duas trajetórias próximas, mas com direções distintas, provocando oscilações em azimute e elevação.
Em 5:21, o sistema passa por um \textit{wrap} do azimute, quando o algoritmo precisa lidar com uma transição de $360^{\circ}$ para $0^{\circ}$ durante a passagem do KSM1-C.

É possível ver o sistema divergindo quando desengajado durante o rastreio do ASRTU-1 (AO-123) até o momento que este entra em sua rotina para retornar à volta inicial simulando o desenrolar dos fios do sistema em 6:57.
O satélite STAR ONE C2 é rastreado em 7:00, sendo este um satélite geoestacionário.
Em 8:28, novamente o \textit{wrap} do azimute é demonstrado ao desengajar um sistema que estava próximo de $0^{\circ}$ (METOP-C) e rastreando um satélite próximo de $360^{\circ}$ (SCD 2).
Em 8:45, é possível perceber o que ocorre quando há uma interrupção no rastreio em uma região onde a trajetória sofre grande aceleração (por estar próxima ao zênite), mostrando que o sistem diverge rapidamente da rota original, tanto em azimute quanto em elevação.
O sistema realiza mais um retorno à posição inicial em 9:39 e permanece por mais tempo sem receber dados, posicionado com valores de azimute e elevação igual a zero e em 10:08 o sistema é completamente desligado das informações do aplicativo até o vídeo terminar.

\end{apendicesenv}
% ---